In questo capitolo vogliamo introdurre il concetto di $k$-superfici e di $k$-forme.

L'esempio più calzante di questi concetti è dato dallo studio dell'elettromagnetismo 
in fisica: i campi vettoriali (come il campo elettrico) si rappresentano tramite $2$-forme 
mentre i campi pseudo-vettoriali (come il campo magnetico) si rappresentano tramite $1$-forme.
Di un campo elettrico vogliamo poter fare il flusso attraverso una superficie $2$-dimensionale.
Di un campo magnetico vogliamo poter fare la circuitazione lungo una curva ($1$-superficie).
La carica elettrica si rappresenta tramite una $3$-forma che vogliamo 
poter integrare su un volume (ovvero una $3$-superficie).
L'obiettivo finale di questo capitolo sarà quello di dimostrare i teoremi di Stokes 
e Gauss-Green, che si possono scrivere in questo modo
\[
  \iint_S (\rot \vec B)\cdot d\vec S 
  = \oint_{\partial S} \vec B \cdot d\vec r, \qquad 
  \iiint_V \div \vec E\, dV 
  = \iint_{\partial V} \vec E\cdot d\vec S.
\]

Ovviamente questi concetti sono fondamentali non solo in fisica ma anche in moltissimi
ambiti della matematica. 
In particolare sebbene il caso delle forme nello spazio tridimensionale saranno 
il nostro obiettivo principale, vogliamo dare delle definizioni 
che si generalizzano ad ogni dimensione spaziale, per questo parleremo di $k$-forme,
e $k$-superfici in $\RR^n$ con $k$ intero, $0\le k \le n$.

Non vogliamo qui fare la teoria delle varietà differenziabili 
(argomento della geometria differenziale) né, tantomeno, 
quella delle correnti rettificabili. 
Per mantenere un livello elementare e per poter dimostrare i teoremi a cui 
stiamo puntando, studieremo una sottoclasse delle correnti rettificabili dato 
dalle superficie parametrizzate regolari a tratti.

L'esempio più importante di superficie che vogliamo essere in grado di rappresentare 
è la sfera. 
Intuitivamente il flusso di un campo vettoriale $\xi$ attraverso una sfera $S$ è l'integrale,
che si può indicare con $\int_S \xi\cdot \nu_S\, d\sigma$,
svolto sulla superficie della sfera, del prodotto scalare tra il campo vettoriale $\xi$
e il versore normale della sfera $\nu_S$.
Il simbolo $d\sigma$ rappresenta la misura di superficie, ovvero l'elemento di area infinitesimo.
La superficie della sfera può essere orientata in due modi opposti se si sceglie 
la normale esterna oppure la normale interna. 
L'integrale del campo vettoriale cambierà segno a seconda dell'orientazione.
Inoltre ci aspettiamo che tale integrale sia lineare rispetto al campo vettoriale:
\[
\int_S (c \xi) \cdot \nu\, d\sigma = c \int_S \xi \cdot \nu\, d\sigma,
\qquad
\int_S (\xi+\eta)\cdot \nu \, d\sigma = \int_S \xi\cdot \nu \, d\sigma + \int_S \eta\cdot\nu\, d\sigma.
\]

La sfera unitaria (superficie bidimensionale nello spazio tridimensionale) può essere rappresentata in molti modi diversi. 
Il modo più banale è quello di considerare l'insieme $\ENCLOSE{(x,y,z)\colon x^2+y^2+z^2=1}\subset \RR^3$.
Purtroppo questa rappresentazione non è ideale per il nostro scopo perché non è in grado di rappresentare 
l'orientazione e la molteplicità. 
In particolare la possibilità di avere una molteplicità è rilevante per avere una struttura algebrica sulle 
superfici: una curva che si avvolge più volte su se stessa (come le spire di un trasformatore elettrico) 
ha una chiara interpretazione fisica che verrebbe persa perché l'unione insiemistica $S \cup S = S$
non distingue due copie sovrapposte della stessa superficie. 
Un problema simile si avrebbe con superfici (o curve) che si autointersecano: 
nei punti di intersezione non sarà possibile definire univocamente un vettore normale alla superficie.

Dal punto di vista matematico quello che si perde, nell'approccio puramente insiemistico, è la struttura 
algebrica delle superfici. Vogliamo infatti permettere la definizione di somma di due superifici $S+T$ 
in modo tale che l'integrale risulti additivo anche rispetto alla superficie. 
Vorremmo inoltre definire un operatore di bordo $\partial S$ che associ ad una $k$-superficie una $(k-1)$-superficie
e che sia tale che $\partial(S+T) = \partial S + \partial T$.

Tutto questo con l'obiettivo finale di ottenere una dimostrazione della formula di Stokes:
\[
 \int_{\partial S} \omega = \int_S d \omega.
\]
 
Denotiamo con $\current \phi$ la \emph{superficie orientata} 
parametrizzata da $\phi$. 



Ad esempio $\phi(x,y) = (x,y,\sqrt{1-u^2-v^2})$ rappresenta 
la semicirconferenza di equazione $x^2+y^2+z^2=1$ con $z\ge 0$.
Un'altra parametrizzazione della stessa semicirconferenza è data, 
in coordinate sferiche, da $\psi(u,v) = ()

